\documentclass{article}
\usepackage[utf8]{inputenc}
\usepackage{booktabs}
\usepackage{longtable}
\title{AAP live one-paper fallback}
\begin{document}
\maketitle
\begin{abstract}
AAP live one-paper fallback synthesizes 1 papers across 5 evidence rows using factor structure as the main organizing lens. The synthesis is concentrated in Survey data from 260 respondents in Dalian, China, behavioral factors, with recurring methods such as Structural equation modeling (SEM) / Survey-based research and task coverage centered on Predicting pedestrian violating crossing behavior intention. Across 1 papers, the evidence is most coherently organized around factor structure rather than any single method, and the strongest clusters are Survey data from 260 respondents in Dalian, China, behavioral factors, Structural equation modeling (SEM) / Survey-based research. Future work should prioritize standardized reporting and more explicit cross-study comparison.
\end{abstract}
\noindent\textbf{Keywords:} predicting pedestrian violating crossing behavior intention, structural equation modeling (sem) / survey-based research, survey data from 260 respondents in dalian, china, examining pedestrians' self-reported violating crossing behavior intentions by applying the theory of planned behavior (tpb) and extending it with descriptive norm, perceived risk, and conformity tendency to evaluate their respective impacts on behavioral intentions., cmin/df ratio, rmsea, pclose, behavioral factors
\section{Introduction and Review Scope}
Introduction and Review Scope provides the analytic entry point for this subsection, bringing the discussion into focus around Behavioral Factors, Survey Data From 260 Respondents In Dalian, China. Across 1 evidence unit, the relevant literature assembles method families such as Structural Equation Modeling (Sem) / Survey Based Research, tasks such as Predicting Pedestrian Violating Crossing Behavior Intention, datasets such as Survey Data From 260 Respondents In Dalian, China, which defines the comparison set for the discussion. Read this way, the subsection begins from a shared problem space rather than a sequence of isolated paper summaries. \cite{08_2016_zhou_aap}

Across the reviewed studies, a recurring pattern emerges in introduction and review scope. Conformity tendency (the tendency to modify behavior when around others) is a strong predictor of violating crossing intention, indicating that the presence of other pedestrians influences behavioral intention. Subjective norm is significant in the basic TPB model but becomes non-significant when descriptive norm is added, suggesting that observation of others' behavior has greater influence than perceived approval from important referents. Collectively, these observations support a subsection-level claim instead of leaving the evidence as a descriptive inventory. \cite{08_2016_zhou_aap}

Taken together, the literature on introduction and review scope points to a field-level pattern rather than a single-study result. The evidence clusters most clearly around Behavioral Factors, Survey Data From 260 Respondents In Dalian, China. The synthesis therefore clarifies both the main takeaway and the boundary conditions that qualify it. \cite{08_2016_zhou_aap}

\section{Factor Taxonomy and Contextual Drivers}
Factor Taxonomy and Contextual Drivers provides the analytic entry point for this subsection, bringing the discussion into focus around Behavioral Factors, Survey Data From 260 Respondents In Dalian, China. Across 1 evidence unit, the relevant literature assembles method families such as Structural Equation Modeling (Sem) / Survey Based Research, tasks such as Predicting Pedestrian Violating Crossing Behavior Intention, datasets such as Survey Data From 260 Respondents In Dalian, China, which defines the comparison set for the discussion. Read this way, the subsection begins from a shared problem space rather than a sequence of isolated paper summaries. \cite{08_2016_zhou_aap}

Across the reviewed studies, a recurring pattern emerges in factor taxonomy and contextual drivers. Conformity tendency (the tendency to modify behavior when around others) is a strong predictor of violating crossing intention, indicating that the presence of other pedestrians influences behavioral intention. Instrumental attitude is a significant predictor of pedestrians' violating crossing behavior intention across all models tested. Collectively, these observations support a subsection-level claim instead of leaving the evidence as a descriptive inventory. \cite{08_2016_zhou_aap}

Comparison is most informative where work on factor taxonomy and contextual drivers converges in broad outline but diverges in operational detail. behavioral factors appears in 5 evidence unit(s), with 0 gap signal(s) and 0 contradiction signal(s). Survey data from 260 respondents in Dalian, China appears in 5 evidence unit(s), with 0 gap signal(s) and 0 contradiction signal(s). The resulting contrast shows where convergence is substantive and where apparent agreement rests on non-equivalent designs. \cite{08_2016_zhou_aap}

Taken together, the literature on factor taxonomy and contextual drivers points to a field-level pattern rather than a single-study result. The evidence clusters most clearly around Behavioral Factors, Survey Data From 260 Respondents In Dalian, China. The synthesis therefore clarifies both the main takeaway and the boundary conditions that qualify it. \cite{08_2016_zhou_aap}

\section{Factor Focus: Behavioral Factors}
Behavioral Factors provides the analytic entry point for this subsection, bringing the discussion into focus around Behavioral Factors, Survey Data From 260 Respondents In Dalian, China. Across 1 evidence unit, the relevant literature assembles method families such as Structural Equation Modeling (Sem) / Survey Based Research, tasks such as Predicting Pedestrian Violating Crossing Behavior Intention, datasets such as Survey Data From 260 Respondents In Dalian, China, which defines the comparison set for the discussion. Read this way, the subsection begins from a shared problem space rather than a sequence of isolated paper summaries. \cite{08_2016_zhou_aap}

Across the reviewed studies, a recurring pattern emerges in behavioral factors. Instrumental attitude is a significant predictor of pedestrians' violating crossing behavior intention across all models tested. Descriptive norm (perceived actual behavior of family and friends) is a stronger predictor than subjective norm (perceived approval) and remains significant when added to the basic TPB model. Collectively, these observations support a subsection-level claim instead of leaving the evidence as a descriptive inventory. \cite{08_2016_zhou_aap}

Taken together, the literature on behavioral factors points to a field-level pattern rather than a single-study result. The evidence clusters most clearly around Behavioral Factors, Survey Data From 260 Respondents In Dalian, China. The synthesis therefore clarifies both the main takeaway and the boundary conditions that qualify it. \cite{08_2016_zhou_aap}

\section{Comparative Factor Evidence and Interactions}
Comparative Factor Evidence and Interactions provides the analytic entry point for this subsection, bringing the discussion into focus around Behavioral Factors, Survey Data From 260 Respondents In Dalian, China. Across 1 evidence unit, the relevant literature assembles method families such as Structural Equation Modeling (Sem) / Survey Based Research, tasks such as Predicting Pedestrian Violating Crossing Behavior Intention, datasets such as Survey Data From 260 Respondents In Dalian, China, which defines the comparison set for the discussion. Read this way, the subsection begins from a shared problem space rather than a sequence of isolated paper summaries. \cite{08_2016_zhou_aap}

Comparison is most informative where work on comparative factor evidence and interactions converges in broad outline but diverges in operational detail. behavioral factors appears in 5 evidence unit(s), with 0 gap signal(s) and 0 contradiction signal(s). Survey data from 260 respondents in Dalian, China appears in 5 evidence unit(s), with 0 gap signal(s) and 0 contradiction signal(s). The resulting contrast shows where convergence is substantive and where apparent agreement rests on non-equivalent designs. \cite{08_2016_zhou_aap}

Across the reviewed studies, a recurring pattern emerges in comparative factor evidence and interactions. Instrumental attitude is a significant predictor of pedestrians' violating crossing behavior intention across all models tested. Descriptive norm (perceived actual behavior of family and friends) is a stronger predictor than subjective norm (perceived approval) and remains significant when added to the basic TPB model. Collectively, these observations support a subsection-level claim instead of leaving the evidence as a descriptive inventory. \cite{08_2016_zhou_aap}

Taken together, the literature on comparative factor evidence and interactions points to a field-level pattern rather than a single-study result. The evidence clusters most clearly around Behavioral Factors, Survey Data From 260 Respondents In Dalian, China. The synthesis therefore clarifies both the main takeaway and the boundary conditions that qualify it. \cite{08_2016_zhou_aap}

\section{Conclusion}
The field-level picture across 1 papers is that factor structure provide the clearest organizing lens, with the strongest signals concentrated in Survey data from 260 respondents in Dalian, China, behavioral factors, Structural equation modeling (SEM) / Survey-based research. \cite{08_2016_zhou_aap}

Stable conclusions: Across 1 papers, the evidence is most coherently organized around factor structure rather than any single method, and the strongest clusters are Survey data from 260 respondents in Dalian, China, behavioral factors, Structural equation modeling (SEM) / Survey-based research. Methodological diversity is real, but the review still shows recurring families such as Structural equation modeling (SEM) / Survey-based research, which suggests the field is comparing similar problems with different operational choices. Reporting remains uneven across tasks and metrics such as Predicting pedestrian violating crossing behavior intention and CMIN/df ratio, RMSEA, which limits direct comparison across studies. \cite{08_2016_zhou_aap}

Unresolved tensions: A second tension is that cross-study comparison still depends on inconsistent reporting conventions. \cite{08_2016_zhou_aap}

Research priorities: Future work should prioritize standardized reporting and more explicit cross-study comparison. A third priority is to align methods, metrics, and context descriptors so the literature can be synthesized more defensibly. \cite{08_2016_zhou_aap}

\appendix
\section{Appendix}
\subsection{Appendix Summary}
\noindent The appendix records 1 papers and 5 matrix rows used to generate the review synthesis. The most visible methodological families are Structural equation modeling (SEM) / Survey-based research, while recurring tasks include Predicting pedestrian violating crossing behavior intention.

\begin{itemize}
\item \textbf{Papers}: 1
\item \textbf{Matrix rows}: 5
\item \textbf{Verified gaps}: 
\item \textbf{Dominant axis}: factor
\item \textbf{Top methods}: Structural equation modeling (SEM) / Survey-based research
\item \textbf{Top tasks}: Predicting pedestrian violating crossing behavior intention
\item \textbf{Top datasets}: Survey data from 260 respondents in Dalian, China
\end{itemize}

\subsection{Evidence Table}
\begin{longtable}{p{0.14\textwidth}p{0.21\textwidth}p{0.10\textwidth}p{0.17\textwidth}p{0.17\textwidth}p{0.17\textwidth}}
\toprule
Paper & Title & Year & Methods & Tasks & Gap matches \\
\midrule
\endfirsthead
\toprule
Paper & Title & Year & Methods & Tasks & Gap matches \\
\midrule
\endhead
08\_2016\_zhou\_aap & An extension of the theory of planned behavior to predict pedestrians' violating crossing behavior using structural equation modeling & 2016 & Structural equation modeling (SEM) / Survey-based research & Predicting pedestrian violating crossing behavior intention & -- \\
\bottomrule
\end{longtable}
\end{document}
% BIB START
@article{08_2016_zhou_aap,
  author={Hongmei Zhou; Stephanie Ballon Romero; Xiao Qin},
  title={An extension of the theory of planned behavior to predict pedestrians' violating crossing behavior using structural equation modeling},
  year={2016},
  journal={Accident Analysis and Prevention}
}